\documentclass[border=8pt]{standalone}
\usepackage{xcolor}
\usepackage{tikz}
\usepackage{enumitem}
\usetikzlibrary{arrows.meta,calc}

% ---------- Farben ----------
\definecolor{codebg}{RGB}{250,250,248}
\definecolor{codekw}{RGB}{60,90,190}
\definecolor{highlightcol}{RGB}{230,80,60}

\begin{document}
\begin{tikzpicture}
    \node[align=left, inner sep=10pt, text width=9.2cm]
          (callout)
          {\textcolor{highlightcol}{\ttfamily\bfseries\footnotesize Kernel-Fenster \& Indexing}\\[6pt]
          {\footnotesize
          \begin{itemize}[leftmargin=1.1em, itemsep=4pt, topsep=0pt]
            \item Die Schleifen laufen über das Kernel-Fenster: \texttt{radius} Pixel in jede Richtung um den aktuellen Pixel \texttt{(x,y)}.
            \item \texttt{(y+y\_)*width + (x+x\_)}: rechnet die 2D-Koordinate in den 1D-Buffer-Index um (row-major).
          \end{itemize}}
          \vspace{10pt}

          \begin{tikzpicture}[x=0.42cm, y=0.42cm]
            % ---------- 11x7 Pixelraster ----------
            \foreach \n in {0,...,7} {
                \draw[black!25, line width=0.4pt] (0,\n) -- (11,\n);
            }
            \foreach \n in {0,...,11} {
                \draw[black!25, line width=0.4pt] (\n,0) -- (\n,7);
            }

            % Buffer-Index (hellgrau) in allen Zellen (row-major: r*width + c)
            % unterste Zeile ausgenommen, da dort der -radius/radius Pfeil liegt
            \foreach \r in {0,...,5} {
                \foreach \c in {0,...,10} {
                    \pgfmathtruncatemacro{\idx}{\r*11+\c}
                    \node[black!45, font=\tiny] at (\c+0.5,6.5-\r) {\idx};
                }
            }

            % Zentrum-Pixel (x,y)
            \fill[codekw, opacity=0.85] (5,3) rectangle (6,4);

            % Kernel-Fenster radius=2 -> 5x5 um das Zentrum
            \draw[highlightcol, line width=1.1pt] (3,1) rectangle (8,6);

            % -radius / radius entlang der Mittelzeile des Fensters
            \draw[{Latex[length=1.6mm]}-{Latex[length=1.6mm]}, black!70, line width=0.6pt]
                (3,0.55) -- (8,0.55);
            \node[anchor=north, font=\tiny] at (3,0.4) {$-$radius};
            \node[anchor=north, font=\tiny] at (8,0.4) {radius};
            \node[anchor=north, font=\tiny] at (5.5,0.4) {$0$};
          \end{tikzpicture}
          };
\end{tikzpicture}
\end{document}